\documentclass[12pt]{article}
\usepackage[T1]{fontenc}
\usepackage{lmodern}

\usepackage[margin=1in]{geometry}
\usepackage{amsmath,amssymb,amsthm,mathtools}
\usepackage{booktabs,array,enumitem,graphicx,microtype,longtable}
\usepackage[colorlinks=true,allcolors=blue]{hyperref}

\ifdefined\pdftrailerid
  \pdftrailerid{[<A8C6D7E83B4F19A7D05C2E164DB8F271>
                 <A8C6D7E83B4F19A7D05C2E164DB8F271>]}
\fi

\newtheorem{theorem}{Theorem}[section]
\newtheorem{proposition}[theorem]{Proposition}
\newtheorem{lemma}[theorem]{Lemma}
\newtheorem{corollary}[theorem]{Corollary}
\theoremstyle{definition}
\newtheorem{definition}[theorem]{Definition}
\theoremstyle{remark}
\newtheorem{remark}[theorem]{Remark}

\newcommand{\QQ}{\mathbb Q}
\newcommand{\rank}{\operatorname{rank}}
\newcommand{\Span}{\operatorname{span}}
\newcommand{\target}{t}

\title{Exact Rank Obstruction and Dictionary-Relative Completion\\
in a Finite Quadratic Response Model}
\author{Luqman Mahmood\thanks{Copyright 2026 Luqman Mahmood. This paper is licensed under Creative Commons Attribution 4.0 International (CC BY 4.0); see \url{https://creativecommons.org/licenses/by/4.0/}. Software is excluded from this grant.}}
\date{1 October 2026}
\hypersetup{pdftitle={Exact Rank Obstruction and Dictionary-Relative Completion in a Finite Quadratic Response Model},pdfauthor={Luqman Mahmood}}

\begin{document}
\maketitle

\begin{abstract}
We resolve an exact linear-algebra question in a fixed rational response
model with $19$ coordinates and $2{,}106$ evaluation functionals.  Nine native
rows have rank $9$; adjoining a distinguished target raises the rank to $10$,
whereas adjoining two formal-family rows gives rank $11$.  An explicit
coordinate functional separates the target from the native coefficient span,
while an exact identity locates it in the formal-family extension.  A second,
rank-$13$ assembly also misses the target.  Five linear equalities and one
nonvanishing condition classify every rational single-row addition that
completes this assembly.  The seven-element dictionary is a specialization:
exactly $r_5$ and $r_7$ complete it, with both identities in closed form.
A nineteen-texture minor has determinant $-2/9$, so evaluation separates the
entire coefficient space.  The classification is confined to this fixed
response model and does not establish physical availability of its rows.
\end{abstract}

\section{Introduction}\label{sec:perimeter}

Finite quadratic response models lead to a basic reachability problem: given
a collection of response rows, does a distinguished target lie in their span,
and if not, which additional generators remove the deficit?  This note studies
one such problem over $\QQ$.  The data consist of a $19$-coordinate coefficient
space, a pool of $2{,}106$ scalar evaluations, nine native rows, a target, two
formal-family rows, and a separate seven-element completion dictionary.  The
goal is not to infer the data from measurements, but to certify exactly what
these fixed data can and cannot span.

The rows are the stipulated inputs of a finite generator-selection problem,
not observations from an experiment.  Their coordinates distinguish scalar,
gradient, and quadratic response terms in the declared model.  This makes
the example a test of which extra response generators resolve a specified
linear obstruction; it does not derive a constitutive law or assert that a
particular apparatus realizes these rows.  No physical interpretation is
attached to these results.

Three features make the example useful.  First, small rational minors certify
a rank hierarchy $9<10<11$ without numerical conditioning questions.  Second,
an explicit coordinate functional identifies the direction missing from the
native rows.  Third, an explicit annihilator calculation classifies all
rational single-row completions of a larger rank-$13$ assembly.  The original
seven-candidate calculation is a corollary; its successful formulas share a
coefficient forced by the native obstruction.  A full-rank nineteen-texture
certificate makes coefficient-space and texture-response statements
equivalent.  The classification uses elementary quotient-space linear algebra;
its content is the explicit criterion for these stipulated rows.

The measurement language is standard.  Model-relative distinguishability and
operator-system span are discussed by Heinosaari, Mazzarella, and Wolf, while
Scott develops informationally complete operator frames and reconstruction
formulas \cite{HeinosaariMazzarellaWolf2013,Scott2006}.  Exact experimental
design over a finite candidate set is likewise an established optimization
setting \cite{AhipasaogluCipollaGondzio2026}.  Here these ideas are used only
in their elementary finite-dimensional form: a nonzero $d\times d$ response
minor proves that $d$ selected scalar evaluations distinguish the specified
$d$-dimensional row family.

All conclusions are relative to the definitions below.  The continuum names
attached to a few rows are mnemonic labels, not a microscopic derivation or a
claim of physical co-realizability.  Section~\ref{sec:model} defines the model
and the evaluation map; Sections~\ref{sec:ranks}--\ref{sec:completion} prove
the rank, obstruction, and completion statements; the supplement supplies the
exact minors, including a machine-readable full-model certificate and a
bounded independent checker.

\section{The fixed response model}\label{sec:model}

\begin{definition}[Coefficient space]
Let $V=\QQ^{19}$ in the ordered basis
\[
(a_0,a_1,a_2,b_1,b_2,b_3,g_1,g_2,
 r_1,r_2,r_3,r_4,r_5,r_6,r_7,m_0,m_1,m_2,\mu_0).
\]
All omitted coordinates in the row definitions below are zero.
\end{definition}

The nine native rows are
\begin{align*}
e_0&=a_0+\tfrac23 b_2+\tfrac13 g_2,\\
e_1&=a_1+b_1-b_3+\tfrac12 g_1-\tfrac12 g_2,\\
e_2&=a_2+b_1-\tfrac23 b_2+b_3+\tfrac12 g_1+\tfrac16 g_2,\\
\beta_0&=\tfrac13r_1-\tfrac13r_3-\tfrac13r_4-\tfrac13r_5+\tfrac23r_6,\\
\beta_1&=\tfrac12r_2+\tfrac12r_5-r_7,\\
\beta_2&=\tfrac23r_1-\tfrac12r_2+\tfrac13r_3-\tfrac23r_4
          -\tfrac16r_5-\tfrac23r_6+r_7,\\
M_s&=2m_0+\tfrac12m_1+\tfrac12m_2,\qquad
M_t=\mu_0,\qquad \kappa=r_1+r_2.
\end{align*}
The two formal-family rows and target are
\begin{equation}\label{eq:formal-target}
F_A=b_2-r_7,\qquad F_B=-\tfrac14r_2,\qquad
\target=-\tfrac23b_2-\tfrac19r_2+\tfrac23r_7.
\end{equation}

The texture pool is also fixed.  Let $\omega\in\{0,1,2\}$, let $k$ range
over the nine integer triples listed in Supplement~\ref{supp:textures}, and let
the twelve-component field vector
have either one unit entry or two unit entries.  The resulting
$3\cdot9\cdot(12+\binom{12}{2})=2{,}106$ textures are ordered
by $\omega$, then by the displayed order of $k$, then by field: the twelve
single-unit fields in coordinate order, followed by the $66$ two-unit fields
in lexicographic order of their occupied-position pairs.  The exact
contractions defining all nineteen response
coordinates are collected in Supplement~\ref{supp:textures}; no floating-point
arithmetic enters the proof.

Write $\mathcal T$ for this ordered texture pool.  Each texture defines a
linear functional on $V$, and together they define the evaluation map
\[
\mathcal E:V\longrightarrow\QQ^{\mathcal T},\qquad
v\longmapsto\bigl(v(\tau)\bigr)_{\tau\in\mathcal T}.
\]
Thus a row has both a coefficient vector in $V$ and a response vector
$\mathcal E(v)$.  Unless stated otherwise, ranks and residual counts below
refer to response vectors on $\mathcal T$; identities such as
Eq.~\eqref{eq:formal-target} are identities in the coefficient space $V$.
A nonzero response minor certifies coefficientwise independence as well,
because any relation in $V$ would remain a relation after applying
$\mathcal E$.

\begin{lemma}[Full-coefficient separation]\label{lem:full-separation}
The evaluation map $\mathcal E$ has rank $19$ and is injective on $V$.
Nineteen textures from the specified pool already separate every pair of
coefficient rows.  In particular, a linear identity holds on all textures
if and only if it holds coefficientwise in $V$.
\end{lemma}

\begin{proof}
In the ordered coordinate basis of $V$, the nineteen texture columns listed
in Supplement~\ref{supp:full-separation} form a $19\times19$ matrix with
determinant $-2/9$.  Its entries are given by the texture contractions and
supplied as exact rational data.  The minor gives rank at least $19$;
$\dim V=19$ gives the upper bound and injectivity.
\end{proof}

\section{Ranks 9, 10, and 11}\label{sec:ranks}

\begin{theorem}[Exact rank hierarchy]\label{thm:rank-hierarchy}
On the texture pool of Section~\ref{sec:model},
\begin{align*}
\rank\{e_0,e_1,e_2,\beta_0,\beta_1,\beta_2,M_s,M_t,\kappa\}&=9,\\
\rank\{e_0,e_1,e_2,\beta_0,\beta_1,\beta_2,M_s,M_t,\kappa,\target\}&=10,\\
\rank\{e_0,e_1,e_2,\beta_0,\beta_1,\beta_2,M_s,M_t,\kappa,F_A,F_B\}&=11.
\end{align*}
\end{theorem}

\begin{proof}
Supplement~\ref{supp:minors} displays exact square response minors with
determinants
\[
\det R_9=\frac29,\qquad
\det R_{10}=-\frac5{81},\qquad
\det R_{11}=-\frac1{36}.
\]
They give the three lower bounds.  The corresponding families contain exactly
$9$, $10$, and $11$ rows, which gives the matching upper bounds.
\end{proof}

\begin{proposition}[Coordinate obstruction]\label{prop:obstruction}
The functional
\begin{equation}\label{eq:obstruction}
\mathcal L=-\tfrac23a_0^*+\tfrac23a_2^*+b_2^*
\end{equation}
annihilates every native row and satisfies $\mathcal L(\target)=-2/3$.
Consequently $\target$ is not in the native coefficient span in $V$.
\end{proposition}

\begin{proof}
Substitution in the nine displayed native rows gives zero; substitution in
Eq.~\eqref{eq:formal-target} gives $-2/3$.
\end{proof}

\begin{proposition}[Formal-family identity]\label{prop:identity}
Coefficientwise in $V$,
\begin{equation}\label{eq:target-identity}
\target=-\frac23F_A+\frac49F_B.
\end{equation}
\end{proposition}

\begin{proof}
The right side is
$-(2/3)(b_2-r_7)+(4/9)(-r_2/4)$, which is exactly
Eq.~\eqref{eq:formal-target}.
\end{proof}

\begin{remark}[Model-relative number of readings]
Every texture supplies one scalar linear functional.  No set of fewer than
$d$ scalar readings can be injective on a $d$-dimensional response space.  The
displayed nonzero minors show that selected sets of $d=9,10,11$ readings attain
the corresponding ranks for the row families above.  This is a statement only
about the fixed model and texture pool, not a lower bound for arbitrary
experiments, encodings, models, or operator families.
\end{remark}

\section{A thirteen-row named assembly}\label{sec:assembly}

Four coordinate rows are appended to the native family:
\[
S_{r_4}=r_4,\qquad H_{b_1}=b_1,\qquad
H_{b_3}=b_3,\qquad H_{\mathrm{dilat}}=b_2.
\]
Let $\mathcal A$ denote the resulting thirteen-row assembly.  It is a second
span problem, not a revision of the native rank-$9$ family.

The labels recall familiar continuum ingredients: parity-even quadratic
torsion models \cite{HayashiShirafuji1979,BlixtHohmannPfeifer2019},
Oseen--Frank and saddle-splay terms \cite{Frank1958,Napoli2010},
Ericksen--Leslie couplings \cite{Leslie1968,GuoEtAl2024}, and generalized
continua \cite{Mindlin1964,GreenRivlin1964}.  In this note they serve only as
mnemonic names for the four finite rows displayed above.  No argument depends
on a microscopic model, material interpretation, or physical co-realization.

We use one fixed residual convention in the completion calculation.  With the
textures in the specified order, exact left-to-right pivot reduction of a
response vector against an ordered row family produces a remainder in
$\QQ^{\mathcal T}$; its \emph{residual count} is the number of nonzero
coordinates.  In the table below, the generator residual reduces
$\mathcal E(r_j)$ against $\mathcal E(\mathcal A)$, while the target residual
reduces $\mathcal E(\target)$ against the assembly augmented by $r_j$.

\begin{proposition}[Rank-$13$ assembly]\label{prop:rank13}
The named assembly $\mathcal A$ has rank $13$.  The target is not in its
span and leaves $792$ nonzero texture residuals after exact pivot reduction.
\end{proposition}

\begin{proof}
The exact $13\times13$ minor in Supplement~\ref{supp:assembly-minors} has
determinant $-1/9$.  The family contains thirteen rows, giving the upper bound.  Exact
pivot back-substitution of the target followed by a scan of all $2{,}106$
responses gives $792$ nonzero residuals.
\end{proof}

\section{Single-row completion in the fixed model}\label{sec:all-completions}

Write $A=\Span(\mathcal A)$ and $B=A+\Span\{\target\}$.  For a coefficient
row $x\in V$, define the five linear forms
\begin{equation}\label{eq:completion-annihilators}
\begin{aligned}
\phi_1(x)&=x_{g_1}-\tfrac12(x_{a_1}+x_{a_2}),\\
\phi_2(x)&=x_{g_2}-\tfrac13x_{a_0}+\tfrac12x_{a_1}-\tfrac16x_{a_2},\\
\phi_3(x)&=2x_{r_3}+x_{r_6},\\
\phi_4(x)&=x_{m_1}-\tfrac14x_{m_0},\\
\phi_5(x)&=x_{m_2}-\tfrac14x_{m_0},
\end{aligned}
\end{equation}
and the additional functional
\begin{equation}\label{eq:assembly-functional}
s(x)=-2x_{r_3}+2x_{r_5}+x_{r_7}.
\end{equation}

\begin{theorem}[All rational single-row completions]\label{thm:all-completions}
For any $x\in V$,
\[
\target\in A+\Span\{x\}
\quad\Longleftrightarrow\quad
\phi_1(x)=\cdots=\phi_5(x)=0\quad\hbox{and}\quad s(x)\ne0.
\]
Equivalently, the completing rows are exactly $B\setminus A$.  For every such
row the coefficient of $x$ is forced, and
\begin{equation}\label{eq:general-completion}
\target-\frac{2}{3s(x)}x\in A.
\end{equation}
The same criterion holds for the corresponding texture-response vectors.
\end{theorem}

\begin{proof}
Substitution shows that all five $\phi_i$ vanish on the thirteen rows of
$\mathcal A$ and on $\target$.  They are independent: their coefficients at
$g_1,g_2,r_6,m_1,m_2$, respectively, form the $5\times5$ identity matrix.
Their common kernel therefore has dimension $14$.  Substitution also gives
$s|_A=0$ and $s(\target)=2/3$.  By Proposition~\ref{prop:rank13},
$\dim A=13$, so $\dim B=14$.  Since $s(\target)\ne0$, the subspace
$B\cap\ker s$ is a hyperplane of $B$ containing $A$, hence equals $A$.
Consequently,
\[
B=\bigcap_{i=1}^5\ker\phi_i,\qquad A=B\cap\ker s.
\]
If the stated conditions hold, $x\in B$ and the left side of
Eq.~\eqref{eq:general-completion} lies in $B\cap\ker s=A$.
Conversely, if $\target=a+cx$ with $a\in A$, then $c\ne0$ because
$\target\notin A$.  Thus $x\in B$ and applying $s$ gives
$2/3=c\,s(x)$, proving both the conditions and the forced coefficient.
Lemma~\ref{lem:full-separation} transfers the equivalence to response vectors.
\end{proof}

Both parts of the criterion are needed.  The row $r_1$ satisfies the five
equalities but has $s(r_1)=0$; the row $r_3$ has $s(r_3)=-2\ne0$ but
$\phi_3(r_3)=2$.  Neither completes the target.  The theorem classifies
algebraically available rows of $V$, not physically realizable generators.

\section{Seven-candidate completion}\label{sec:completion}

The candidate dictionary is the seven-element set
$\mathcal D=\{r_1,\ldots,r_7\}$.  Each candidate is appended singly to
$\mathcal A$.  The following specialization of
Theorem~\ref{thm:all-completions} retains the dictionary-specific identities
and the complete residual enumeration.

\begin{corollary}[Dictionary-relative completion]\label{thm:completion}
Exactly $r_5$ and $r_7$ complete the named assembly to the target.  Their
augmented ranks are $14$, witnessed by determinants $-2/9$ and $-1/9$.
The exact identities are
\begin{align}
\target={}&-\tfrac29\beta_0-\tfrac89\beta_1-\tfrac29\beta_2
 +\tfrac29\kappa-\tfrac29S_{r_4}-\tfrac23H_{\mathrm{dilat}}
 +\tfrac13r_5,\label{eq:r5-recipe}\\
\target={}&\phantom{-}\tfrac19\beta_0+\tfrac19\beta_1+\tfrac19\beta_2
 -\tfrac19\kappa+\tfrac19S_{r_4}-\tfrac23H_{\mathrm{dilat}}
 +\tfrac23r_7.\label{eq:r7-recipe}
\end{align}
\end{corollary}

\begin{proof}
The rows $r_1,r_2,r_4$ fail the nonzero test in
Theorem~\ref{thm:all-completions}; $r_3,r_6$ fail its third equality.
The rows $r_5,r_7$ satisfy all five equalities and have $s$-values $2,1$,
respectively, so they are exactly the completing dictionary members.
For each $r_j$, exact pivot reduction determines whether the appended row
raises rank and whether the target residual then vanishes on all $2{,}106$
textures.  The complete table is
\begin{center}
\begin{tabular}{@{}ccccc@{}}
\toprule
candidate & augmented rank & generator residual & target residual & closes \\
\midrule
$r_1$ & 13 & 0   & 792 & no \\
$r_2$ & 13 & 0   & 792 & no \\
$r_3$ & 14 & 234 & 756 & no \\
$r_4$ & 13 & 0   & 792 & no \\
$r_5$ & 14 & 792 & 0   & yes \\
$r_6$ & 14 & 819 & 234 & no \\
$r_7$ & 14 & 792 & 0   & yes \\
\bottomrule
\end{tabular}
\end{center}
Supplement~\ref{supp:assembly-minors} displays both nonzero rank-$14$ minors.
Back-substitution gives Eqs.~\eqref{eq:r5-recipe} and \eqref{eq:r7-recipe};
direct substitution kills all nineteen coefficient residuals and the full
texture scan kills all $2{,}106$ response residuals.
\end{proof}

\begin{remark}[The obstructed direction]\label{rem:forced-coefficient}
The functional $\mathcal L$ from Eq.~\eqref{eq:obstruction} annihilates the
nine native rows, $S_{r_4}$, $H_{b_1}$, $H_{b_3}$, and every candidate
$r_j$, while $\mathcal L(H_{\mathrm{dilat}})=1$.  Hence any coefficient-space
expansion of $\target$ using $\mathcal A$ and one candidate from $\mathcal D$
must give
$H_{\mathrm{dilat}}$ the coefficient
$\mathcal L(\target)=-2/3$.  The common coefficient in
Eqs.~\eqref{eq:r5-recipe} and \eqref{eq:r7-recipe} is therefore forced by the
obstruction, rather than selected by the enumeration.  Subtracting the two
identities also shows that $r_5-2r_7\in\Span(\mathcal A)$: the successful
candidates determine proportional nonzero directions in the quotient
$V/\Span(\mathcal A)$.
This $-2/3$ coefficient law is dictionary-relative, not a statement for every
$x\in V$.  For example, the allowed algebraic choice $x=\target$ completes
with no assembly contribution.  In a general expansion $\target=a+cx$, the
coefficient of $H_{\mathrm{dilat}}$ in $a$ is
$-2/3-c\mathcal L(x)$; it equals $-2/3$ for the dictionary because
$\mathcal L(r_j)=0$.
\end{remark}

\section{Conclusion}\label{sec:limits}

The fixed response model has a simple exact structure.  Its nine native rows
have rank $9$, the target adds one independent direction, and two formal rows
contain that target through an explicit coefficient identity.  A larger
rank-$13$ assembly still misses the target.  Five independent annihilator
conditions and one nonzero test classify every rational single-row completion
within $V$.  Precisely two members of the seven-element dictionary pass.
The native obstruction functional explains the coefficient shared by those
two formulas, while the general criterion identifies the required quotient
direction.  Nineteen specified textures separate the entire coefficient
space, so this classification is equivalent in coefficient and response form.

All calculations use rational arithmetic.  An independent implementation
reconstructed the texture pool, response rows, ranks, obstruction values,
residual counts, and completion identities directly from the definitions.
The supplement records the rank certificates and provides the full-model
minor as exact rational data.  The accompanying bounded checker independently
reconstructs that minor from the four-index projector formulas, verifies the
annihilator criterion and dictionary identities, and tests negative controls.
It does not replay the complete $2{,}106$-texture residual counts.

The result is deliberately finite and model-relative.  The classification
covers all rational single-row additions in this $19$-coordinate model, but
does not classify generators in larger response spaces or arbitrary
multi-row additions.  It proves neither universal measurement minimality,
stability under perturbation, nor a physical realization of the mnemonic
continuum ingredients.  Extension to larger response models requires their
own row definitions and separation certificates.

\appendix
\section{Exact response and certificate supplement}\label{supp:exact}

This supplement records the texture contractions and exact square minors used
in the proofs, including the accompanying machine-readable full-model minor.
Rows and columns occur in the orders stated in the text.  All texture indices
are zero-based.

\setcounter{MaxMatrixCols}{14}
\newcommand{\certmatrix}[1]{%
\begin{center}
\resizebox{\linewidth}{!}{$\begin{pmatrix}#1\end{pmatrix}$}
\end{center}}

\subsection{Texture contractions}\label{supp:textures}

Order the field vector as
\[
(w^0_0,w^0_1,w^0_2,w_{00},w_{01},w_{02},w_{10},w_{11},w_{12},
 w_{20},w_{21},w_{22}).
\]
Let $\omega\in\{0,1,2\}$ and
\[
k\in\{(0,0,0),(1,0,0),(0,1,0),(0,0,1),(1,1,0),(1,0,1),
(0,1,1),(1,1,1),(1,2,3)\}.
\]
The field vector has either one unit entry or two unit entries, yielding
$78$ fields and $2{,}106$ textures.  The order is $\omega$ first, then $k$
in the displayed order, then the twelve single-unit fields in coordinate
order followed by the $66$ pairs in lexicographic order of their occupied
positions.  This is not lexicographic sorting of the resulting binary
field vectors.  Define
\[
\tau=\sum_a w_{aa},\quad |k|^2=\sum_i k_i^2,\quad
d_a=\sum_i k_iw_{ai},\quad \ell_i=\sum_a k_aw_{ai},\quad
k\mathbin{\cdot}w^0=\sum_a k_aw^0_a.
\]
With $P_0,P_1,P_2$ the trace, antisymmetric, and symmetric-traceless
projectors,
\begin{align*}
(P_0)_{ai,bj}&=\delta_{ai}\delta_{bj}/3,\\
(P_1)_{ai,bj}&=(\delta_{ab}\delta_{ij}-\delta_{aj}\delta_{ib})/2,\\
(P_2)_{ai,bj}&=(\delta_{ab}\delta_{ij}+\delta_{aj}\delta_{ib})/2
                -\delta_{ai}\delta_{bj}/3.
\end{align*}
The nineteen scalar responses are
\begin{align*}
a_s&=\omega^2\langle w,P_sw\rangle &(s&=0,1,2),\\
b_1&=\omega\sum_{ai}w_{ai}k_iw^0_a,
&b_2&=\omega\tau(k\mathbin{\cdot}w^0),\\
b_3&=\omega\sum_{ai}w_{ai}k_aw^0_i,
&g_1&=|k|^2|w^0|^2,\\
g_2&=(k\mathbin{\cdot}w^0)^2,
&r_1&=-|k|^2\sum_{ai}w_{ai}^2,\\
r_2&=-|k|^2\tau^2,
&r_3&=-|k|^2\sum_{ai}w_{ai}w_{ia},\\
r_4&=-\sum_a d_a^2,
&r_5&=-\sum_i\ell_i^2,\\
r_6&=-\sum_a d_a\ell_a,
&r_7&=-\left(\sum_{ai}k_ak_iw_{ai}\right)\tau,\\
m_s&=-\langle w,P_sw\rangle &(s&=0,1,2),\\
\mu_0&=|w^0|^2.&&
\end{align*}

\subsection{Full-coefficient separation certificate}\label{supp:full-separation}

Let $u_1,\ldots,u_{19}$ be the ordered coordinate basis of $V$ from
Section~\ref{sec:model}.  Select the following texture indices, in this order:
\[
\begin{gathered}
J=(0,3,4,45,78,79,81,82,84,85,86,123,149,\\
705,706,747,794,798,805).
\end{gathered}
\]
If $j_m$ is the $m$-th entry of $J$, define
\[
M_{\ell m}=\mathcal E(u_\ell)_{j_m},\qquad 1\leq\ell,m\leq19.
\]
Every entry is obtained from the contractions in
Supplement~\ref{supp:textures}.  Exact rational elimination gives
\[
\det M=-\frac29.
\]
The file \path{data/full_model_separating_minor.csv}, relative to the package
root, supplies all $361$ entries.  Its header records the ordered indices
$J$; its first column names the coordinate rows.  Integers and fractions
are exact values, not rounded decimal measurements.  The displayed formulas
and these data specify every entry of the certificate.

The script \path{check_response.py} reconstructs all nineteen columns
independently, using rational arithmetic and direct four-index projector
contractions, compares every CSV entry, and checks rank $19$ and determinant
$-2/9$.  It also verifies the five independent annihilators, $s|_A=0$,
$s(\target)=2/3$, and the dictionary specialization.  Its negative controls
include a row with $s\ne0$ that violates an annihilator, a row in the assembly span with
$s=0$, a deficient texture selection, and corrupted certificate data.
The checker is bounded: it does not replay the full residual enumeration or
all six smaller displayed minors.  It uses only the Python standard library;
the README gives the command and the output boundary.

\subsection{Exact rank 9, 10, and 11 minors}\label{supp:minors}

Column integer $j$ denotes texture index $j$ in the specified zero-based
ordered construction.  All entries below are rational numbers.

\paragraph{Rank 9.}
The rows are
\[
e_0,e_1,e_2,\beta_0,\beta_1,\beta_2,\qquad M_s,M_t,\kappa.
\]
The columns are $0,3,78,79,81,82,86,149,747$.
\certmatrix{
0&0&1/3&0&0&0&0&0&4/3\\
0&0&0&1/2&0&0&0&0&0\\
0&0&2/3&1/2&0&0&0&0&2/3\\
0&0&0&0&0&0&-1/3&0&0\\
0&0&0&0&0&-1/2&0&-2&0\\
0&0&0&0&0&-1/2&-2/3&0&0\\
0&-1&0&0&-1&-1/2&-1/2&-3&-3\\
1&0&1&1&0&0&0&0&0\\
0&0&0&0&-2&-1&-1&-6&0
}
Its determinant is $2/9$.

\paragraph{Rank 10.}
Append the target and use columns $0,3,78,79,81,82,84,86,149,747$.
\certmatrix{
0&0&1/3&0&0&0&0&0&0&4/3\\
0&0&0&1/2&0&0&0&0&0&0\\
0&0&2/3&1/2&0&0&0&0&0&2/3\\
0&0&0&0&0&0&0&-1/3&0&0\\
0&0&0&0&0&-1/2&0&0&-2&0\\
0&0&0&0&0&-1/2&0&-2/3&0&0\\
0&-1&0&0&-1&-1/2&-1/2&-1/2&-3&-3\\
1&0&1&1&0&0&0&0&0&0\\
0&0&0&0&-2&-1&-1&-1&-6&0\\
0&0&0&0&-5/9&0&0&0&4/9&0
}
Its determinant is $-5/81$.

\paragraph{Rank 11.}
Append $F_A,F_B$ and use columns $0,3,78,79,81,82,84,85,86,149,747$.
\certmatrix{
0&0&1/3&0&0&0&0&0&0&0&4/3\\
0&0&0&1/2&0&0&0&0&0&0&0\\
0&0&2/3&1/2&0&0&0&0&0&0&2/3\\
0&0&0&0&0&0&0&0&-1/3&0&0\\
0&0&0&0&0&-1/2&0&-1/2&0&-2&0\\
0&0&0&0&0&-1/2&0&-1/2&-2/3&0&0\\
0&-1&0&0&-1&-1/2&-1/2&-1&-1/2&-3&-3\\
1&0&1&1&0&0&0&0&0&0&0\\
0&0&0&0&-2&-1&-1&-2&-1&-6&0\\
0&0&0&0&1&0&0&0&0&0&0\\
0&0&0&0&1/4&0&0&1/4&0&1&0
}
Its determinant is $-1/36$.

\subsection{Rank 13 and completion minors}\label{supp:assembly-minors}\label{supp:candidates}

The rank-$13$ rows are
$e_0,e_1,e_2,\beta_0,\beta_1,\beta_2,M_s,M_t,\kappa,
S_{r_4},H_{b_1},H_{b_3},H_{\mathrm{dilat}}$ and the columns are
$0,3,78,79,81,82,84,86,149,747,794,798,805$.
\certmatrix{
0&0&1/3&0&0&0&0&0&0&4/3&4/3&4/3&0\\
0&0&0&1/2&0&0&0&0&0&0&0&0&0\\
0&0&2/3&1/2&0&0&0&0&0&2/3&8/3&2/3&2\\
0&0&0&0&0&0&0&-1/3&0&0&0&0&0\\
0&0&0&0&0&-1/2&0&0&-2&0&0&-1/2&-1/2\\
0&0&0&0&0&-1/2&0&-2/3&0&0&0&-1/2&-1/2\\
0&-1&0&0&-1&-1/2&-1/2&-1/2&-3&-3&-1&-1&-1/2\\
1&0&1&1&0&0&0&0&0&0&1&1&1\\
0&0&0&0&-2&-1&-1&-1&-6&0&-2&-2&-1\\
0&0&0&0&-1&0&-1&0&0&0&-1&0&0\\
0&0&0&0&0&0&0&0&0&0&1&0&0\\
0&0&0&0&0&0&0&0&0&0&1&0&1\\
0&0&0&0&0&0&0&0&0&0&1&1&0
}
Its determinant is $-1/9$.

For $r_5$, append its response row and use the columns
\[
0,3,78,79,81,82,84,85,86,149,747,794,798,805.
\]
\certmatrix{
0&0&1/3&0&0&0&0&0&0&0&4/3&4/3&4/3&0\\
0&0&0&1/2&0&0&0&0&0&0&0&0&0&0\\
0&0&2/3&1/2&0&0&0&0&0&0&2/3&8/3&2/3&2\\
0&0&0&0&0&0&0&0&-1/3&0&0&0&0&0\\
0&0&0&0&0&-1/2&0&-1/2&0&-2&0&0&-1/2&-1/2\\
0&0&0&0&0&-1/2&0&-1/2&-2/3&0&0&0&-1/2&-1/2\\
0&-1&0&0&-1&-1/2&-1/2&-1&-1/2&-3&-3&-1&-1&-1/2\\
1&0&1&1&0&0&0&0&0&0&0&1&1&1\\
0&0&0&0&-2&-1&-1&-2&-1&-6&0&-2&-2&-1\\
0&0&0&0&-1&0&-1&0&0&0&0&-1&0&0\\
0&0&0&0&0&0&0&0&0&0&0&1&0&0\\
0&0&0&0&0&0&0&0&0&0&0&1&0&1\\
0&0&0&0&0&0&0&0&0&0&0&1&1&0\\
0&0&0&0&-1&-1&0&0&0&0&0&-1&0&-1
}
Its determinant is $-2/9$.

For $r_7$, use the same columns and append the $r_7$ response row.
\certmatrix{
0&0&1/3&0&0&0&0&0&0&0&4/3&4/3&4/3&0\\
0&0&0&1/2&0&0&0&0&0&0&0&0&0&0\\
0&0&2/3&1/2&0&0&0&0&0&0&2/3&8/3&2/3&2\\
0&0&0&0&0&0&0&0&-1/3&0&0&0&0&0\\
0&0&0&0&0&-1/2&0&-1/2&0&-2&0&0&-1/2&-1/2\\
0&0&0&0&0&-1/2&0&-1/2&-2/3&0&0&0&-1/2&-1/2\\
0&-1&0&0&-1&-1/2&-1/2&-1&-1/2&-3&-3&-1&-1&-1/2\\
1&0&1&1&0&0&0&0&0&0&0&1&1&1\\
0&0&0&0&-2&-1&-1&-2&-1&-6&0&-2&-2&-1\\
0&0&0&0&-1&0&-1&0&0&0&0&-1&0&0\\
0&0&0&0&0&0&0&0&0&0&0&1&0&0\\
0&0&0&0&0&0&0&0&0&0&0&1&0&1\\
0&0&0&0&0&0&0&0&0&0&0&1&1&0\\
0&0&0&0&-1&0&0&0&0&0&0&-1&0&0
}
Its determinant is $-1/9$.


\section*{Acknowledgment of assistance}
AI-assisted research, computation, checking, and editorial support used
OpenAI Codex and Anthropic Claude. Luqman Mahmood directed the work
and is responsible for the claims and final text. These internal
checks are not external peer review.

\bibliographystyle{plain}
\bibliography{references}

\end{document}
